\chapter{Dialysis}

\section*{What Is This Test or Procedure?}
Dialysis is the process of removing excess water, solutes, and toxins from the blood in patients whose kidneys can no longer perform these functions naturally (kidney failure).

\section*{Alternative Names}
Renal replacement therapy, Hemodialysis, Peritoneal dialysis

\section*{Principle}
Blood is cleansed across a semipermeable membrane: by diffusion against dialysis fluid in hemodialysis, or across the patient's own peritoneal membrane in peritoneal dialysis. Waste products and excess fluid move out of the blood and electrolytes are rebalanced.
\section*{History}
Willem Kolff constructed the first working dialyzer in 1943 during the German occupation of the Netherlands, saving the first patient's life in 1945.

\section*{Why This Test or Procedure Is Ordered}
Ordered for acute kidney injury or end-stage renal disease (ESRD) to manage severe uremia, hyperkalemia, volume overload, or metabolic acidosis.

\section*{Is This a Primary or Secondary Test or Procedure}
Primary treatment for life-threatening uremia, severe fluid overload, or drug toxicities.

\section*{How This Test or Procedure Is Performed}
Hemodialysis pumps blood through an external filter (dialyzer) and returns it. Peritoneal dialysis uses the lining of the abdomen as a filter, inserting dialysate fluid via a catheter.

\section*{What Preparation Is Needed}
A vascular access (AV fistula, graft, or temporary catheter) must be surgically created prior to starting hemodialysis.

\section*{Contraindications and Precautions}
No absolute contraindications for life-saving dialysis. Peritoneal dialysis is contraindicated in patients with severe abdominal adhesions.

\section*{Risks}
Hypotension, muscle cramps, infection of the access site, bleeding, electrolyte imbalances, or dialysis disequilibrium syndrome.

\section*{What Happens After the Test or Procedure}
Monitor blood pressure, weight (to assess fluid removal), and post-dialysis laboratory values (BUN, creatinine, electrolytes).

\section*{Understanding Your Results}
Efficacy is measured by Kt/V (a measure of urea clearance). A Kt/V of at least 1.2 is target for hemodialysis.

\section*{Normal Results}
Adequate reduction in blood urea nitrogen (BUN) and creatinine; correction of electrolyte abnormalities (potassium, bicarbonate); normal volume status.

\section*{Follow-up Tests and Procedures}
Regular nephrology clinic visits, access site monitoring, and dietary/fluid management.

\section*{References}
\begin{enumerate}
  \item MedlinePlus. Dialysis. U.S. National Library of Medicine; 2025.
  \item Current Medical Diagnosis and Treatment. McGraw-Hill; 2025.
\end{enumerate}

\clearpage
